\section*{Chapter \ref{ch:b2:externalities-and-public-goods-theory} --- Externalities and Public Goods: The Formal Theory}

\begin{solution}{exo:b2:externalities-and-public-goods-theory:1}
% figdata: so2_summary
From 14.76 to 0.79 million tonnes, a fall of 95 per cent. Other things changed over the same years: other rules on the
same pollutant, the fuels burned and the technology of power plants; the chart has no comparison group showing what
emissions would have been without the programme.
\end{solution}

\begin{solution}{exo:b2:externalities-and-public-goods-theory:2}
% code: exercises/solutions.py externality
Unregulated: $100 - 2E = 0$, $E = 50$. Efficient: $100 - 2E = 2E$, $E = 25$. Tax $2 \times 25 = 50$. Loss $\tfrac12(2 + 2)(50 - 25)^2 = 1{,}250$.
\end{solution}

\begin{solution}{exo:b2:externalities-and-public-goods-theory:3}
Lighthouse: public (no ship's use reduces another's, and passing ships cannot be kept from seeing it). Music service: club
(non-rival, but excludable by subscription). Open fishing ground: common-pool (rival, since fish caught by one are not
caught by another, but not excludable). Sandwich: private (rival and excludable).
\end{solution}

\begin{solution}{exo:b2:externalities-and-public-goods-theory:4}
% code: exercises/solutions.py pivot
Net values 150, 60 and $-60$: the bridge is built. For each resident, the others' net values sum to 0, 90 and 210, all at least zero:
no one is pivotal and no one pays a tax. With a cost of 330, shares are 110 and net values 140, $-50$ and $-70$, summing to 20: built.
Without resident 1 the others sum to $-120$, so she is pivotal and pays 120, keeping $140 - 120 = 20$; the others are not pivotal.
\end{solution}

\begin{solution}{exo:b2:externalities-and-public-goods-theory:5}
The mill's gain from 0 to 40 tonnes is $\int_0^{40}(120 - 2E)\,dE = 4{,}800 - 1{,}600 = 3{,}200$; the fishery's damage is $\int_0^{40}E\,dE = 800$. Any payment from the mill
to the fishery between 800 and 3{,}200 makes both better off.
\end{solution}

\begin{solution}{exo:b2:externalities-and-public-goods-theory:6}
% code: exercises/solutions.py permits_two_firms
Abatement from 60 to $E$ costs firm A $(60 - E_A)^2$ and firm B $(60 - E_B)^2/2$. With 25 each: $35^2 + 35^2/2 = 1{,}225 + 612.5 = 1{,}837.5$. With trading,
marginal costs are equal: $120 - 2E_A = 60 - E_B$ with $E_A + E_B = 50$, so $E_A = 110/3 = 36.67$, $E_B = 13.33$, and the price is $120 - 2 \times 36.67 = 46.67$. The
cost is $23.33^2 + 46.67^2/2 = 1{,}633.3$, a saving of 204.2. Firm A, which abates expensively, buys permits from firm B.
\end{solution}

\begin{solution}{exo:b2:externalities-and-public-goods-theory:7}
% code: exercises/solutions.py samuelson
The sum is $36 - 4G = 12$, so $G^* = 6$; \omterm{def:b2:externalities-and-public-goods-theory:lindahl}{personalised prices} $12 - 6 = 6$, 4, 2 and 0. Voluntarily, household 1 buys until $12 - G = 12$: none, and
no one else buys either. The loss is $\int_0^6 (24 - 4G)\,dG = 144 - 72 = 72$.
\end{solution}

\begin{solution}{exo:b2:externalities-and-public-goods-theory:8}
% code: exercises/solutions.py weitzman
For the average case $\bar E = 90/9 = 10$ and the tax is 60. Efficient levels: $105/9 = 11.67$ and $75/9 = 8.33$. Under the tax the mill emits
$(105 - 60)/3 = 15$ or 5. Losses: tax $\tfrac12 \times 9 \times (10/3)^2 = 50$ in each case, since $15 - 35/3 = 10/3$; quota $\tfrac12 \times 9 \times (5/3)^2 = 12.5$. The quota is better, as the
slopes rule says ($d = 6 > b = 3$; ratio $36/9 = 4$).
\end{solution}

\begin{solution}{exo:b2:externalities-and-public-goods-theory:9}
% ledger: b2:cheung1973
For cherries the hive gives the grower pollination and the beekeeper almost no honey, so the grower pays; for mint it gives
the beekeeper honey and the grower nothing, so the beekeeper pays. Among neighbours, the spillover from each orchard's
hives reaches several others; contracts among all of them would cost more to negotiate than a custom of matching hive
numbers that each grower can see, enforced by the inconveniences neighbours impose on a ``bad neighbour''.
\end{solution}

\begin{solution}{exo:b2:externalities-and-public-goods-theory:10}
% code: exercises/solutions.py identical_gap
Household $i$ buys more only while $a_i - G > c$. In equilibrium $G$ satisfies $a_i - G \le c$ for all $i$, with equality for anyone who buys. The
household with the largest $a_i$ buys until $G = a_1 - c$; at that $G$ every other has $a_i - G < c$ (or equal, buying nothing more), so $G = \max(a_1 - c, 0)$.
With $n$ identical households, $G^* = a - c/n$ and the voluntary $G = a - c$: the gap $c(1 - 1/n)$ rises with $n$. With $a = 10$, $c = 6$, $n = 3$: 8 against 4.
\end{solution}

\begin{solution}{exo:b2:externalities-and-public-goods-theory:11}
% code: exercises/solutions.py river
It sets $90 - 3E = 3E$: a quota of 15 tonnes. The loss is $\tfrac12(3 + 1.5)(20 - 15)^2 = 56.25$, a quarter of the 225 lost without regulation.
\end{solution}

\begin{solution}{exo:b2:externalities-and-public-goods-theory:12}
With $\mathit{MD} = d$ constant (a flat line), the efficient level satisfies $\mathit{MB}(E) = d$ for every $\theta$; a tax of $d$ makes the polluter choose exactly
that. In the formula, the tax's loss has $d^2$ for the slope of $\mathit{MD}$, which is zero. If $\mathit{MD}$ is vertical at a threshold $\bar E$, the efficient level is
$\bar E$ whatever $\theta$ (as long as $\mathit{MB}(\bar E)$ is positive), and a quota at $\bar E$ reaches it; in the formula, the quota's loss $\theta^2/(2(b + d))$ goes to zero as
$d$ grows. A tax suits damage that accumulates slowly, a quantity one with a threshold.
\end{solution}

\begin{solution}{pb:b2:externalities-and-public-goods-theory:1}
% code: exercises/solutions.py river
\textbf{Part I.}
1. $90 - 3E = 0$: 30 tonnes.
2. $\int_0^{30}(90 - 3E)\,dE = 2{,}700 - 1{,}350 = 1{,}350$ thousand dalens a month.
3. $\int_0^{30}1.5E\,dE = 675$.
4. $1{,}350 - 675 = 675$.
5. The fishery's damage: it reaches the fishery through the river, not through any price the mill pays.
\textbf{Part II.}
6. $90 - 3E = 1.5E$: 20 tonnes.
7. Benefit $1{,}800 - 600 = 1{,}200$, damage 300, net 900.
8. $900 - 675 = 225$.
9. The mill loses $\int_{20}^{30}(90 - 3E)\,dE = 150$; the fishery gains $\int_{20}^{30}1.5E\,dE = 375$. Any payment from the fishery between 150 and 375.
10. The mill gains 1{,}200 from 20 tonnes and the fishery loses 300: any payment from the mill between 300 and 1{,}200.
11. In both, the joint surplus $\int_0^E (90 - 4.5E)\,dE$ is largest at 20, and any other level leaves a gain both could share.
12. Many fishers to agree with; the mill's costs and the fishery's damages are private; agreements are hard to enforce.
\textbf{Part III.}
13. $1.5 \times 20 = 30$.
14. $30 \times 20 = 600$.
15. $1{,}200 - 600 = 600$.
16. 30, the mill's marginal benefit at 20 tonnes.
17. The discharge is still 20; the mill keeps the permits' value, 600, instead of paying it.
18. $90 - 3E = 3E$: 15 tonnes; the loss is $\tfrac12 \times 4.5 \times 25 = 56.25$.
\textbf{Part IV.}
19. $(90 \pm 15)/4.5$: 23.33 and 16.67.
20. $(90 \pm 15 - 30)/3$: 25 and 15.
21. $\tfrac12 \times 4.5 \times (25 - 70/3)^2 = \tfrac12 \times 4.5 \times (5/3)^2 = 6.25$ in each case.
22. $\tfrac12 \times 4.5 \times (10/3)^2 = 25$ in each case.
23. The tax: 6.25 against 25. Slopes rule: ratio $d^2/b^2 = 2.25/9 = 0.25$.
24. With $4.5E$: average efficient level $90/7.5 = 12$, tax 54. The tax's loss is $225 \times 20.25/(2 \times 9 \times 7.5) = 33.75$, the quota's $225/15 = 15$: the
quota is better, as $d = 4.5 > b = 3$.
25. Efficient discharge 20 tonnes a month, against 30 unregulated, with 225 thousand dalens a month lost without regulation;
a tax of 30 or 20 permits reaches it; with uncertain costs the tax loses 6.25 in expectation and the quota 25, because damage
rises more slowly than the mill's marginal benefit falls.
\end{solution}
